\section{实验结果}

% Page 1: Experimental Setup
\begin{frame}{计算实验}
\tiny 数据来自DB1B\cite{DOT_DB1B_2025}和AOTP\cite{DOT_OTP_2025}，经过合成得到 Frontier Airlines 的实例。
\begin{columns}[T,totalwidth=\textwidth]
\column{0.55\textwidth}
\begin{block}{Frontier Airlines (F9)}
\begin{itemize}
    \item 美国超低成本航空公司（ULCC）
    \item 60 个机场，259 条航段
    \item 382 个 OD 市场 （起点-终点对）
    \item Airbus A320 系列机队（4 个子类型）
    \item 座位容量 186--228 座
    \item 最多 1 次中转，最大 3 小时中转时间
\end{itemize}
\end{block}

\column{0.38\textwidth}
\begin{block}{实例规模}
\centering
\small
\begin{tabular}{lcc}
\toprule
实例 & 时间跨度 & 候选行程 \\
\midrule
% F9-0.5x & 12h & 81,252 \\
F9-1.0x & 24h & 475,776 \\
F9-1.5x & 36h & 557,028 \\
F9-2.0x & 48h & 951,552 \\
F9-2.5x & 60h & 1,032,804 \\
F9-3.0x & 72h & 1,427,328 \\
\bottomrule
\end{tabular}
\vspace{0.1cm}
\end{block}

\vspace{0.1cm}
\tiny
离散化粒度为 15 分钟\\
所有实例均在相同硬件配置进行\\
AMD EPYC 9654 96-Core Processor，512GB RAM
\end{columns}
\end{frame}


% Page 2: Main results table
\begin{frame}{计算性能对比}
\tiny
Raw Model：使用 Gurobi 直接求解原始模型\\
IA-DR：使用我们提出的 IA-DR 框架求解\\
Gap 定义为 $|\frac{\text{UB}-\text{LB}}{\text{LB}}|$
\begin{block}{Raw Model vs. IA-DR}
\centering
\small
\begin{tabular}{lcccc}
\toprule
实例 & Raw Gap (2h) & IA-DR Gap (2h) & Raw Gap (12h) & IA-DR Gap (12h) \\
\midrule
% F9-0.5x & 9.44\%  & 11.79\% & 3.10\%  & 6.11\%  \\
F9-1.0x & 18.40\% & 1.84\%  & 4.75\%  & $<$1\% (2.1h) \\
F9-1.5x & 13.09\% & $<$1\% (0.1h) & $<$1\% (5.7h) & $<$1\% (0.1h) \\
F9-2.0x & --      & 1.25\%  & 15.26\% & 1.25\%  \\
F9-2.5x & 19.65\% & 1.91\%  & 18.06\% & 1.89\%  \\
F9-3.0x & --      & $<$1\% (1.3h) & -- & $<$1\% (1.3h) \\
\bottomrule
\end{tabular}
\vspace{0.1cm}

\footnotesize\tiny ``--'' 表示 Raw Model 在时限内未产生可用的可行解或界。
\end{block}

\vspace{0.15cm}
\small
\begin{itemize}
    \item 实例越大，IA-DR 优势越明显
    \item F9-3.0x（142万 itineraries）：Raw Model 无法求解，IA-DR 在 1.3h 内达到 $<$1\%
\end{itemize}
\end{frame}


% Page 3: B&B tree comparison
\begin{frame}{F9-1.0x分支定界树对比（Gurobi \& IA-DR上界模型）}
\begin{columns}[T,totalwidth=\textwidth]
\column{0.55\textwidth}
\begin{center}
    \includegraphics[width=\linewidth]{../figure/img/unexpl_raw_end_both.pdf}
\end{center}

\column{0.42\textwidth}
\begin{block}{观察}
\begin{itemize}
    \item Raw Model 的搜索树持续膨胀
    \item IA-DR 上界在搜索树早期收敛到下界模型提供的紧下界
\end{itemize}
\end{block}

% \vspace{0.2cm}
% \begin{block}{关键机制}
% \begin{itemize}
%     \item Raw Model：弱 LP 松弛 $\Rightarrow$ 大量无效分支
%     \item IA-DR：紧下界 $\Rightarrow$ 搜索空间大幅缩小
% \end{itemize}
% \end{block}
\end{columns}
\end{frame}



